% fig-ii-11.tex — II.11: cut a line in extreme and mean ratio (golden section).
\begin{figure}[H]
\centering
\begin{tikzpicture}[scale=1.2, line cap=round, line join=round]
  \coordinate (A) at (0, 0);
  \coordinate (B) at (3, 0);
  \coordinate (C) at (0, 3);
  \coordinate (D) at (3, 3);
  % Square ABDC on AB.
  \draw[very thick] (A) -- (B) -- (D) -- (C) -- cycle;
  % E = midpoint of AC.
  \coordinate (E) at (0, 1.5);
  \draw[thin, dashed] (E) -- (B);
  % F = on AC produced, with EF = EB. EF = sqrt(9 + 2.25) ~ 3.354.
  \coordinate (F) at (0, {-sqrt(11.25) + 1.5});
  \draw[thin] (E) -- (F);
  % H on AB with AH = AF.
  \coordinate (H) at ({sqrt(11.25) - 1.5}, 0);
  \draw[very thick, dotted] (H) -- ($(H)+(0, 3)$);
  % Labels.
  \node[above left]  at (A) {$A$};
  \node[above right] at (B) {$B$};
  \node[below right] at (D) {$D$};
  \node[below left]  at (C) {$C$};
  \node[left]        at (E) {$E$};
  \node[left]        at (F) {$F$};
  \node[below]       at (H) {$H$};
  % Arc from F to H (centre A, radius AF).
  \draw[thin] (F) arc[start angle=270, end angle=360, radius={sqrt(11.25) - 1.5}];
\end{tikzpicture}
\caption{Proposition II.11. Square $ABDC$ on $AB$; midpoint $E$ of $AC$;
$F$ on the extension of $AC$ with $EF = EB$. Then $AH = AF$ cuts $AB$
in the desired ratio: $AB \cdot HB = AH^2$. This is the golden section.}
\label{fig:II.11}
\end{figure}
